% Independently audited by path_stability and random_hosts.
% Exact certificates: four-regular-saturated-small-cut-checks.json.
\section{A saturated terminal pair produces a small cyclic cut}
\label{sec:saturated-pair-small-cut}

\begin{theorem}[A full pair inside a fixed-order pole]
\label{thm:saturated-pair-small-cyclic-cut}
Let $X$ be a connected normalized $k$-pole of rank $h$, with $k\ge2$ distinct ports.
Suppose one terminal pair has $2^h$ simple paths.
If $h>5k$, then every cubic completion of $X$ has a cyclic edge cut of size at most four.
Here a completion adds one incidence at each port and preserves all edges of $X$.
\end{theorem}

\begin{proof}
Keep the saturated pair as terminals and suppress the other $k-2$ ports.
The resulting two-pole $N$ still has rank $h$ and a full terminal-path count.
It has no loop: the independent loop coordinate would halve the maximum possible path density.
Its positive-rank bridgeless pieces therefore form a series of saturated capped ladders and digons.
Write their ranks as $b_1,\ldots,b_m$, so $\sum_i b_i=h$.

We construct vertex-disjoint cyclic candidates inside these pieces.
If a piece has rank $b\le4$, take its whole vertex set as one candidate.
Its boundary has size at most two in every completion, provided no additional port subdivides an internal edge.
Its two attachment vertices account for these two boundary incidences.

Now consider a piece of rank $b\ge5$.
Its ladder has $b-1$ positions, arranged in singleton and square blocks.
Take the four vertices of every square block as one candidate.
Within each maximal run of singleton blocks, pair consecutive blocks and take their four vertices as further candidates.
Each candidate contains a four-cycle and has boundary at most four.
The candidates are vertex-disjoint.

Let $a$ be the number of square blocks.
Let $s$ be the number of singleton positions, and let $o$ count singleton runs of odd length.
The number of candidates is
\[
 q=a+\frac{s-o}{2}.
\]
Since $o\le a+1$ and $a\le q$, one has
\[
 b-1=2a+s=2q+o\le3q+1.
\]
Consequently,
\[
 q\ge\frac{b-2}{3}\ge\frac b5.
\]
For $b\le4$, the single whole-piece candidate also gives at least $b/5$ candidates.
Summing over the pieces produces at least $h/5$ disjoint cyclic candidates.

Restoring the other $k-2$ ports subdivides at most $k-2$ edges of $N$.
Call a candidate spoiled if a restored port subdivides one of its internal edges.
Each subdivided edge spoils at most one candidate, because their vertex sets are disjoint.
Thus at most $k-2$ candidates are spoiled.
Subdivisions of boundary edges do not spoil a candidate or enlarge its boundary.
The new subdivision vertices simply remain outside its selected vertex set.

Since $h/5>k$, at least two unspoiled candidates remain.
Choose one as a vertex set in the completed cubic graph.
It contains a cycle, and its boundary has size at most four.
The other unspoiled candidate supplies a cycle in its complement.
This boundary is therefore a cyclic edge cut.
\end{proof}

\begin{corollary}[A short proof for five-cuts with two large shores]
\label{cor:five-cut-large-shores-saturation-shortcut}
There are absolute constants $H,K$ with the following property.
Let $G$ arise by joining corresponding ports of connected normalized five-poles of ranks $h,k\ge H$.
Put $r=h+k+4$.
Then
\[
 c(G)\le2^{r-1}+Kr2^{r/2}.
\]
If neither shore has a full terminal pair, then $c(G)<2^{r-1}$.
\end{corollary}

\begin{proof}
Choose $H>25$.
If either shore has a full terminal pair, Theorem~\ref{thm:saturated-pair-small-cyclic-cut} supplies a cyclic cut of size at most four in $G$.
Theorem~\ref{thm:sharp-order-four-cut-error} then gives the asserted error bound.

Suppose neither shore has a full terminal pair.
For one shore, normalize all counts by its affine-space size.
Let $S$ be the sum of its ten path coordinates $a_e$.
Let $F$ be the sum of its five four-terminal forest coordinates $U_A$.
Theorem~\ref{thm:fixed-port-aggregate} gives
\[
 S+F\le8+16(h+4)^{-2^{-145}}-2^{-h}=8+o(1),
\]
where $h$ denotes this shore's rank.
Each forest coordinate lies in $[0,1]$, so $F\le5$.
The full-or-$15/16$ path gap gives $a_e\le15/16$ for every pair.
Hence the shore energy satisfies
\[
 E=\sum_e a_e^2+\sum_A U_A^2
 \le\frac{15}{16}S+F
 \le\frac{15}{16}(8+o(1))+\frac5{16}
 =\frac{125}{16}+o(1).
\]

Apply this to both shores.
The exact five-cut gluing formula and Cauchy--Schwarz bound their normalized crossing contribution by the average of their energies.
It is therefore at most $125/16+o(1)$, strictly below the leading threshold $8$.
The internal cycles have normalized contribution at most $2^{-h}+2^{-k}$.
Increase $H$ until the combined errors are smaller than $3/16$.
Then $c(G)<8\,2^{h+k}=2^{r-1}$, as claimed.
\end{proof}

This argument covers two large shores.
It does not remove the separate bounded-shore problem for five-edge cuts.
The stronger profile and terminal-avoidance lemmas remain valid independent refinements.
